\maketitle

\thispagestyle{empty}

\vspace{6.5cm}

\begin{center}
    \zihao{3}
    \songti
    \textbf{V1.0.3 大学物理B} 

    \textbf{课后习题集 | \LaTeX}
\end{center}

\newpage

\thispagestyle{empty}

\begin{flushleft}
	\songti
	\zihao{3}
	\textbf{附件1}
\end{flushleft}

\begin{center}
	\songti
	\zihao{3}
	\textbf{2025-2026年第二学期大学物理B考试大纲}\footnote{画删除线的内容表示期末考试不涉及.}
\end{center}

\begin{table}[H]
	\centering
	\songti
	\renewcommand{\arraystretch}{1.38}
	\begin{tabularx}{\textwidth}{
			>{\centering\arraybackslash}p{0.35\textwidth}
			>{\centering\arraybackslash}X
		}
		\toprule
		\textbf{章节} & \textbf{考试范围} \\
		\midrule
		第1章\quad 质点运动学
		& 习题1-4、1-5、1-12、1-18、1-20、1-25、1-26 \\
		
		第2章\quad 质点动力学
		& 习题2-5、2-8、2-11、2-15、2-17、2-25、2-29、2-31、2-38 \\
		
		第3章\quad 刚体力学基础
		& 习题3-10、3-11、3-12、3-13、3-14、3-16、3-17 \\
		
		第4章\quad 机械振动与机械波
		& 习题4-11、4-15、4-17、4-18、4-19、4-21；
		习题4-7、4-22、4-24、4-25、4-26、4-30 \\
		
		\sout{第5章\quad 气体动理论}
		& \sout{例题5.2、例题5.3、例题5.4；
		习题5-2、5-4、5-5、5-8、5-10} \\
		
		\sout{第6章\quad 热力学基础}
		& \sout{例题6.1、例题6.2、例题6.3；
		习题6-6、6-11、6-15、6-17} \\
		
		第7章\quad 真空中的静电场
		& 习题7-9、7-11、7-13、7-14、7-15、7-17；
		例题7.2、例题7.5、例题7.7、例题7.9、例题7.12、例题7.11 \\
		
		第8章\quad 导体和电介质
		& 例题8.1、例题8.2；
		习题8-5 \\
		
		第9章\quad 恒定磁场
		& 例题9.2；
		习题9-10、9-12、9-15、9-16、9-19、9-20 \\
		
		第10章\quad 变化的电磁场
		& 习题10-20、10-22、10-23；
		例题10.1、例题10.2、例题10.5、例题10.7；
		习题10-8、10-17 \\
		
		第11章\quad 波动光学
		& 例题11.1、例题11.2；
		习题11-27、11-30、11-35；
		例题11.5；
		习题11-43、11-44、11-45、11-47 \\
		\bottomrule
	\end{tabularx}
\end{table}

\newpage

\pagestyle{plain}
\setcounter{page}{1}
\tableofcontents
\newpage
\setcounter{page}{1}
